\section{Vulnerability Replay Evidence and Rediscovery Protocol}
\label{sec:cve}
The supplied \path{Key_Experiment/vulnerability} collection contains four case dossiers: three labeled with CVE identifiers and one SMARTPL availability case without a CVE. Each includes a report and a trigger input; three include vulnerable-side log excerpts, and ProFTPD also includes a patched-side replay log. Tables~\ref{tab:cve} and~\ref{tab:cveresults} report this evidence. Identifiers, versions, and replay counts are attributed to these local records; the archive does not include the upstream patches or complete build manifests needed to independently certify every vulnerability/patch pair.

\subsection{Case identity and evidence provenance}
We distinguish an attached log observation (L) from a replay result stated only in a case report (R). A report can contain additional trials not represented by the supplied log excerpt. Repeated lines within a log are not independent trials. Input hashes and byte lengths are recorded in the companion \path{vulnerability_evidence_20260929.json}; no additional target executions were performed for this manuscript update. Table~\ref{tab:cve} identifies the recorded builds, triggering conditions, and differential evidence.

\begin{table*}[!tbp]
\centering\small
\caption{Vulnerability cases and archived evidence (L: log; R: report).}
\label{tab:cve}
\begin{tabularx}{\textwidth}{@{}>{\raggedright\arraybackslash}p{0.20\textwidth} >{\raggedright\arraybackslash}p{0.19\textwidth} Y Y@{}}
\toprule
Case / interface & Reported vulnerable build & Trigger evidence / oracle & Patch or control evidence \\
\midrule
\mbox{CVE-2023-51713}\newline ProFTPD / FTP & \texttt{61e621e743346};\newline ASan; 65,535-byte command buffer & 60,006-byte input; heap-buffer-overflow: one-byte read beyond a 65,568-byte region (L). & Base + \texttt{97bbe68363cc}; three rejections, no ASan (L); same input/buffer. \\
\addlinespace[3pt]
\mbox{CVE-2026-38998}\newline LIVE555 / RTSP & Campaign: 2023.05.10; public report: 2026.02.26 (R) & 728-byte public PoC; heap-use-after-free (L). Fuzzer-discovered input absent. & Local handler-removal patch: 1/1 still triggers (R); upstream-patch differential absent. \\
\addlinespace[3pt]
\mbox{CVE-2026-39863}\newline Kamailio / SIP over TCP & \texttt{a2209018fb};\newline TCP listener required & 278-byte input; parsed Content-Length $-10$ (L): overflow predicate, not crash count. & Local \texttt{INT\_MAX/10} guard suppresses fingerprint (R); reported fix \texttt{060b6c2530c8}. No patched log/repeat count. \\
\addlinespace[3pt]
SMARTPL availability case\newline forked-daapd / HTTP & Reported 27.2,\newline \texttt{2ca10d9b} & 148-byte input; lexer errors; health probe fails; process alive (R). Replay log absent. & No patched replay/CVE; parser reportedly removed in 28.4. \\
\bottomrule
\end{tabularx}
\end{table*}

For ProFTPD, the report identifies an out-of-bounds read in \texttt{make\_ftp\_cmd}; the attached ASan trace records the corresponding heap-buffer-overflow signature. The oracle uses a 65,535-byte command buffer. The report states that the default 512-byte configuration does not expose the same error to ASan, so the outcome is conditional on this configuration. The patched log lists three rounds returning an FTP 500 response with no ASan report and normal one-shot exit. This supports a recorded differential replay; immutable image digests, build flags, and the actual applied patch remain unavailable in this collection.

For LIVE555, the supplied ASan excerpt contains the public-PoC use-after-free signature at offset \texttt{0x614b86}. Its two-line observation is repeated verbatim and is counted once as an attached signature. The report separately attributes fuzzer-generated inputs to the same allocation/free lifecycle, reports 8/9 replay hits, and records failure of a local patch attempt. The described \path{trigger_loopfuzz.raw}, full allocation/free traces, and per-trial replay metadata are absent. The archived public PoC therefore supports replay diagnostics, but cannot demonstrate independent rediscovery by the evaluated controller or establish the cross-version CVE attribution on its own.

For Kamailio, the attached log records a negative parsed Content-Length, providing evidence for the reported overflow predicate. This path requires TCP and cannot be evaluated through the UDP-only benchmark configuration. The latest report records suppression by a local guard, but supplies neither a patched-side log nor a repetition denominator. That observation is retained as R; it does not verify equivalence to the identified upstream patch. The report's earlier registry-only 20/20 versus 0/20 statement is excluded from the quantitative table because the registry and trial records are not supplied. The SMARTPL report describes reproducible unavailability from a known input, while none of five sampled campaign hang inputs satisfies that predicate. This case has neither a CVE identity nor a verified patched counterpart.

\subsection{Replay outcomes and their denominators}
Table~\ref{tab:cveresults} preserves each reported batch separately. A numerator counts replay attempts or selected queue inputs satisfying that case's predicate; the denominator is the number tested. These are not independent 12-hour fuzzer runs, and the main experiment's nominal $N=10$ is not substituted for the actual replay denominator. Selected batches, per-input repeats, and later aggregate counts may overlap; they are not pooled into one success rate. In particular, 19/20 selected inputs is not 95\% campaign recall, and 8/9 LIVE555 replays is below the prospective 95\% stability threshold.

\begin{table*}[!tbp]
\centering\small
\caption{Reported replay batches and separate controls; original denominators retained.}
\label{tab:cveresults}
\begin{tabularx}{\textwidth}{@{}>{\raggedright\arraybackslash}p{0.20\textwidth} Y Y Y@{}}
\toprule
Case & Vulnerable-side batches (R) & Patched or diagnostic control & Interpretation / missing evidence \\
\midrule
\mbox{CVE-2023-51713} & 1/1 seed; 1/1 fuzzer crash (Sept.~22); 3/3 batch crashes (Sept.~25). & Reported fix: 0/3 ASan activations; 3/3 FTP 500 rejections (L, Sept.~28). & Separate replay batches; no arm-specific discovery probability or trigger time. \\
\addlinespace[3pt]
\mbox{CVE-2026-38998} & 1/1 public PoC; three fuzzer inputs: 3/3 each; later aggregate: 8/9. & Local patch: 1/1 still triggers (R); successful patched-side denominator absent. & Unpooled batches; fuzzer input/trial records absent; no confirmed patch pair. \\
\addlinespace[3pt]
\mbox{CVE-2026-39863} & 1/1 seed; queue samples: 3/4 (60 s), 19/20 (60 min), 19/20 (1,540 min), 50/51 (separate 60-min batch). & Local guard: fingerprint absent (R); repeat count unavailable. & Queue samples, not campaign recall; no extrapolation to 195 entries; patched log absent. \\
\addlinespace[3pt]
SMARTPL availability case & 5/5 advisory repeats; 3/3 48-message-seed repeats; 1/1 later advisory replay. & Sampled campaign hangs: 0/5 SMARTPL hits (R); patched replay absent. & Known-input unavailability; no confirmed campaign rediscovery or CVE. \\
\bottomrule
\end{tabularx}
\end{table*}

The reports also list crash-file totals and campaigns outside the supplied benchmark/ablation ledger. These totals are not unique vulnerabilities, and durations such as 60 or 1,540 minutes do not identify time to first trigger. The collection lacks a join from each reproducing input to a verified A--E run, its model information history, and its independent oracle/root cause. We therefore do not assign these replay outcomes to individual arms, claim newly discovered CVEs, or infer admission/calibration gains from them.

\subsection{Controlled rediscovery still to be measured}
The confirmatory benchmark requires a vulnerable base and the same base plus a minimal upstream security patch, with dependencies, instrumentation, configuration, and input compatibility held constant. Each case must have a network-reachable trigger, an independent reached/triggered oracle, reproducible builds/reset, and patch validation. The planned six-case pilot across at least three protocols, followed by expansion toward twelve, remains unfinished. Pilot qualification requires at least 95\% observed oracle repeatability over a declared validation set, no patched-target activation, automated build/reset/cleanup, and information separation. The four dossiers above are case evidence, not six qualified benchmark pairs; a small all-success replay batch alone does not establish an underlying success probability of at least 95\%.

Information-controlled campaigns must exclude CVE identifiers, advisories, public PoCs, and patch contents from model inputs and campaign seeds, while retaining anonymous oracle IDs in separate benchmark metadata. The public PoC and known triggers in this collection are validation inputs; their replay is not evidence of information-controlled rediscovery. Source-level trigger conditions, sanitizer reports with patch validation, and protocol invariants distinguish reaching relevant code from actually triggering a defect. Experiments require isolated target networks and reset between trials; these are protocol requirements, not properties inferred from the reports.

For each CVE and arm, the formal plan retains ten independent runs and a 12-hour deadline, frozen before execution. Deadline recall, time-to-trigger, reached-but-not-triggered counts, patched-target activation, executions/coverage/tokens before trigger, and recall per CPU-hour or 100,000 tokens remain unavailable for the controlled A--E comparison. Nontriggers require right-censoring records, and interrupted runs require separate failure accounting. The supplied replay batches do not supply those risk sets, so Kaplan--Meier estimates or arm-level security-effectiveness claims are not reported.
